\chapter{Releases}\label{ch:releases}

\emph{You want several issues built one after another, unattended, and shipped together.} A
release is a named plan of issues for one project, in build order. Once started, Agern builds it
on its own branch \texttt{release/<name>}: one run per issue, each merged into the release
branch. When nothing is left to build, Agern opens one pull request from the release branch to
the default branch, and the release is merged when you approve it, or automatically. The rules in
full are in \repodoc{ARCHITECTURE.md\#releases}{\texttt{ARCHITECTURE.md}, ``Releases''}.

\section{Create a release}\label{sec:create-release}

\begin{enumerate}
  \item On the project page, the \ui{Releases} section of the agents \& runs pane lists the releases
    that are not finished, with their state and progress (figure~\ref{fig:project-releases}).
    Choose its \ui{+ New}.
  \item Give the release a \ui{Name}. Letters, digits, \texttt{.}, \texttt{\_} and \texttt{-} are
    allowed; the name becomes the branch \texttt{release/<name>}. Add a \ui{Description} if you like.
  \item Under \ui{When it is ready}, choose the release mode (section~\ref{sec:release-mode}):
    \ui{Manual approval} or \ui{Release automatically}. Choose \ui{Create release}
    (figure~\ref{fig:new-release}).
\end{enumerate}

Creating a release creates no branch and starts nothing. It is a plan, shown as \ui{in preparation},
that anyone who can see the project may change.

\screenshot[0.55\linewidth]{36-project-releases-desktop}{The releases of a project in the agents \& runs pane.}{project-releases}

\screenshot{41-new-release-desktop}{Creating a release.}{new-release}

\section{Assign issues and set the build order}\label{sec:assign}

\begin{enumerate}
  \item Assign issues in either place: in the project's issue list, with the release selector on
    each issue (or select several issues and use \ui{Assign to\ldots}); or on the release page, with
    \ui{Add an open issue} and \ui{Add}. An issue is in at most one release; assigning it again
    moves it.
  \item On the release page, the issues are listed \ui{in build order}
    (figure~\ref{fig:release-planning}). Move an issue up or down with the arrows, or take it out
    with \ui{Remove}. Put issues that others build on first.
\end{enumerate}

\screenshot{37-release-planning-desktop}{A release in preparation: its issues in build order.}{release-planning}

\section{Choose the release mode}\label{sec:release-mode}

The \ui{Release mode} selector on the release page (figure~\ref{fig:release-mode}) decides what
happens when everything is built:

\begin{itemize}
  \item \ui{Manual approval}: the release waits as \ui{ready to release} until whoever started it
    approves it.
  \item \ui{Automatic}: once everything is built, merges cleanly and all checks on the release pull
    request have passed, Agern merges it with the owner's GitHub token. Held issues, conflicts and
    failed checks stop it.
\end{itemize}

While the release is in preparation anyone may change the mode. Once it is started or scheduled,
only its owner may switch it to automatic, because it merges with the owner's token; the owner or
an admin may switch it back to manual.

\screenshot[0.85\linewidth]{40-release-mode-desktop}{The release mode of a release in preparation.}{release-mode}

\section{Start or schedule it}\label{sec:start-release}

\begin{enumerate}
  \item Choose \ui{Start release} and confirm. Agern creates the branch \texttt{release/<name>} from
    the default branch and starts the run of the first issue. You become the release's owner: all
    its runs use your Claude credential and GitHub token.
  \item Or choose \ui{Schedule start\ldots} and pick a time, for example tonight. At that time
    Agern starts the release as you. If another release of the project is still building then, it
    waits for it. \ui{Change time} and \ui{Cancel schedule} change your mind.
\end{enumerate}

Only one release per project builds at a time.

\section{Follow the build, pause, and resolve held issues}\label{sec:release-building}

The release page of a building release shows the branch, how many issues are merged and held, a
status line such as \ui{Building \#58 \ldots}, and each issue's state: \ui{waiting},
\ui{running}, \ui{merged}, \ui{held} or \ui{skipped}, with a link to its run
(figure~\ref{fig:release-building}).

\begin{enumerate}
  \item To stop starting new issues, choose \ui{Pause}; the running issue finishes. \ui{Resume}
    continues. When your Claude usage limit is reached, Agern pauses the release itself and
    resumes it when the limit resets.
  \item An issue whose run needs a human, or failed, is \ui{held}, with the reason, and the release
    goes on with the next issue. Resolve it in one of these ways:
    \begin{itemize}
      \item open its \ui{run} and approve it (section~\ref{sec:needs-human}): it merges into the
        release branch;
      \item when its finished work could only not be merged because the release branch moved on (a
        migration with the same version, or a merge conflict), choose \ui{Update and merge}: in its
        turn, Agern merges the release branch into the run's branch, the builder resolves the
        conflict, the run is reviewed and approved again, and it merges. A building release already
        tries this once by itself before it holds the issue, so the button is there when that try did
        not work out;
      \item choose \ui{Retry} on the issue and pick how (section~\ref{sec:retry});
      \item or choose \ui{Remove} to leave it out of this release.
    \end{itemize}
  \item \ui{Cancel release} stops the release; its branch stays on GitHub.
\end{enumerate}

\screenshot{38-release-building-desktop}{A building release: one issue merged, one running, one held, one waiting.}{release-building}

\section{Approve and release}\label{sec:approve-release}

When nothing is left to build, the release is \ui{ready to release} (figure~\ref{fig:release-ready}).
Agern has opened the \ui{Release pull request} and shows whether it merges cleanly and the state of
its checks (CI on GitHub). A release with held issues is incomplete and cannot be approved yet.

\begin{enumerate}
  \item Check the pull request on GitHub if you like; \ui{Check again} refreshes the merge state
    and checks.
  \item Choose \ui{Approve and release} and confirm. Agern merges the pull request so that each
    issue stays one commit on the default branch, closes the issues and deletes the release branch.
    The release is \ui{released} and moves to the project's release history.
\end{enumerate}

With \ui{Automatic} mode, Agern does this step itself once the checks pass.

\screenshot{39-release-ready-desktop}{A release that is ready: its pull request, merge state, checks and the approve button.}{release-ready}

\subsection{Releases in workspaces}\label{sec:workspace-releases}

In a workspace project a release can mix issues of several members and of the workspace repository
itself. Each issue needs its member label, as for a run: Agern refuses to plan an issue without
one (in a \texttt{.agern/workspace.yml} workspace), with several, or for a member you cannot read,
and says which label to add. The project agent can plan such releases for you in the chat.

Each issue is built on \texttt{release/<name>} in its own repository; Agern creates that branch in a
member when the first of its issues is built. When the release is ready, the release page lists
one pull request per repository under \ui{Pull requests}, with its state and checks, and the
release is ready only when all of them merge cleanly. Approving merges the member pull requests
first, then moves the submodule pointers in the workspace to the members' merge commits and merges
the workspace pull request. If something fails partway, Agern stops, says which repositories were
merged and which were not, and leaves the rest open: fix the problem and approve again, and Agern
continues with what is not merged. Cancelling such a release deletes its branch in every
repository.

\whatyouhave
You can plan a release, order its issues, start or schedule it, keep it moving past held issues,
and release it by hand or automatically. Next, learn to keep an eye on all agents and on your
usage: chapter~\ref{ch:watch}.
